% =============================================================================
% PRÄAMBEL (praeambel.tex)
% =============================================================================

% --- GRUNDEINSTELLUNGEN ---
\usepackage[utf8]{inputenc}
\usepackage[T1]{fontenc}
\usepackage[ngerman]{babel}
\usepackage{graphicx}
\usepackage{csquotes}
\usepackage{datatool}

% --- MATHEMATIK ---
\usepackage{amsmath, amssymb, amsthm, latexsym, derivative}

% --- GRAFIKEN & DIAGRAMME ---
\usepackage{epsfig, tikz, subcaption, caption, eso-pic}
\usepackage[export]{adjustbox} 
\usepackage[section]{placeins}
\usepackage[european]{circuitikz}

% --- SI-EINHEITEN ---
\usepackage{siunitx}
\sisetup{locale = DE}
\DeclareSIUnit{\euro}{\text{\euro}}

% --- TABELLEN ---
\usepackage[table]{xcolor} 
\usepackage{booktabs, multirow, colortbl}
\usepackage{tabularx}
\usepackage{makecell}

% --- LITERATURVERZEICHNIS (BIBLATEX) ---
\usepackage[backend=biber, style=ieee]{biblatex}
\usepackage{xurl}

% --- SONSTIGE PAKETE ---
\usepackage{enumitem}
\usepackage{makeidx}
\usepackage{totpages}
\usepackage{chngcntr}
\usepackage[normalem]{ulem}
\usepackage{calc}
\usepackage{pdfpages}
\usepackage{blindtext}
\usepackage{float}
\usepackage{newfloat}
\usepackage[most]{tcolorbox}
\usepackage[framemethod=TikZ]{mdframed}
\usepackage{etoolbox}
\usepackage{eurosym}
\usepackage{verbatim}

% --- HYPERREF (sollte fast immer als letztes geladen werden) ---
\usepackage[colorlinks=true, linkcolor=blue, citecolor=green, urlcolor=magenta, hidelinks]{hyperref}

% --- GLOSSAR-KONFIGURATION ---
% Wichtig:symbols-Typ definieren
\usepackage[toc, acronym, symbols, nonumberlist, section]{glossaries}


% Eigenen Stil definieren
\newglossarystyle{mynonbold}{
  \setglossarystyle{list}
  \renewcommand*{\glossentryname}[1]{\glsentryname{##1}}
}

\makeglossaries

% --- LAYOUT-DATEI LADEN ---
% (Wird in der Haupt-Datei geladen)

% --- EINSTELLUNG FÜR KOMA-SKRIPT ---
\KOMAoptions{numbers=noenddot}

% --- ABSTÄNDE ---
\setlength{\parindent}{0em}
\setlength{\parskip}{0.5em}

% --- FLOAT-UMGEBUNGEN ---
\DeclareFloatingEnvironment[listname={Diagrammverzeichnis}, name=Diagramm, fileext=lod, placement=htp]{diagram}
\DeclareFloatingEnvironment[listname={Oszillogrammverzeichnis}, name=Oszillogramm, fileext=loo, placement=htp]{oszillo}

% --- ZÄHLER-EINSTELLUNGEN ---
\counterwithin{figure}{section}
\counterwithin{table}{section}
\counterwithin{oszillo}{section}
\counterwithin{diagram}{section}

% --- BOX-DEFINITIONEN ---
\newmdenv[linecolor=red, linewidth=1pt, frametitle=Hinweis, frametitlebackgroundcolor=red, frametitlefont=\color{white}\bfseries, frametitlerule=true, backgroundcolor=white, roundcorner=2pt]{hinweisbox}
\newmdenv[linecolor=gray, linewidth=1pt, frametitle=Aufgabe, frametitlebackgroundcolor=white, frametitlefont=\color{black}\bfseries, frametitlerule=true, backgroundcolor=white, roundcorner=2pt]{aufgabenbox}

\newcounter{rechnungbox}[section]
\renewcommand{\therechnungbox}{\thesection.\arabic{rechnungbox}}
\newtcolorbox{rechnungbox}[1][]{colframe=darkgray, colback=white, coltitle=white, colbacktitle=gray, title={Rechnung \therechnungbox}, fonttitle=\bfseries, boxed title style={size=small}, boxrule=1pt, sharp corners=south, before upper=\stepcounter{rechnungbox}, #1}

% --- TOGGLES ---
\newtoggle{showtoc} \toggletrue{showtoc}
\newtoggle{showtables} \toggletrue{showtables}
\newtoggle{showfigures} \toggletrue{showfigures}
\newtoggle{showoszillos} \togglefalse{showoszillos}
\newtoggle{showdiagrams} \togglefalse{showdiagrams}

% --- SCHRIFTART AKTIVIEREN ---
\AtBeginDocument{\haupttextSchrift}

% --- SUBSSUBSECTION DEFINITION ---
\makeatletter
\newcommand{\subsubsubsection}{\@startsection{paragraph}{4}{\z@}%
    {-3.25ex \@plus -1ex \@minus -.2ex}%
    {1.5ex \@plus .2ex}%
    {\normalfont\normalsize\bfseries}}
\makeatother

\setcounter{secnumdepth}{4}
\setcounter{tocdepth}{4}

% --- MATH COMMANDS ---
\newcommand*\diff{\mathop{}\!\mathrm{d}}
\newcommand*\Diff[1]{\mathop{}\!\mathrm{d^#1}}